\begin{abstract}
    This thesis addresses the fundamental incompatibility of applying gradient
    projection-based continual learning (CL) methods to hypernetwork architectures,
    establishing a novel framework for stable subspace optimization.  Traditional 
    projection methods collapse under hypernetwork architectures: the architectural 
    necessity of weight-chunking forces the model to accumulate an exceptionally large 
    gradient volume before a single projection step can execute, blowing up the condition
    number of the projection matrix and causing catastrophic inversion failures that 
    standard regularisation cannot rescue.
    
    The primary research objective is to eliminate these numerical singularities and
    systematically evaluate whether combining hard geometric projections on top of
    task-conditioned hypernetworks and their output-scoped functional regularizers
    yields a constructive synergy.  To achieve this, we introduce
    \emph{projection-scoped normalization}: gradient columns stored in the memory
    matrix $G$ are rescaled to unit $\ell_2$ norm, and the diagonal empirical Fisher
    vector is rescaled by its absolute maximum, together stabilizing the metric tensor
    within the local projection algebra.  We also introduce \texttt{iFOPNG}, an inertial
    variant of Fisher-Orthogonal Projected Natural Gradient Descent \citep{garg_fisher-orthogonal_2026}
    that replaces the current-task Fisher $F_{\text{new}}$ with the combined curvature
    $F_c = F_{\text{new}}+F_{\text{old}}$, inspired by the parameter-inertia concept
    of Elastic Weight Consolidation.  Methods are benchmarked on sequential Split and Permuted MNIST
    , Split-CIFAR10 and Split-CIFAR100. The empirical findings reveal that projection normalization
    resolves matrix inversion collapse in the tested configurations, and that \texttt{iFOPNG} consistently 
    outperforms \texttt{FOPNG} across all benchmarks; in the standard-capacity hypernetwork
    setting only \texttt{iFOPNG} surpasses plain \texttt{SGD}, while most other projection 
    variants fall below it.  A standalone hypernetwork optimized
    with \texttt{Adam} and \citeauthor{oswald_continual_2020}'s functional regularizer
    consistently surpasses all projection variants, indicating that momentum, which
    projection methods forgo by design, is the decisive factor in compressed generative
    weight spaces.
\end{abstract}

